\section{Tracking the actual metric flow}\label{sec:tracking}

\begin{proof}[Proof of Theorem~\ref{thm:main}]
The construction in Section~\ref{sec:approximation} defines
$\widehat H_t$ from the global finite trajectory. It gives uniform
matrix bounds on each bounded slow-time interval and proves
\eqref{eq:approx-thick-intro}.
Lemma~\ref{lem:existence-bounds} gives the actual solution for all
physical times and preserves its initial bounds $cI\le H_t\le CI$.
Thus actual and comparison metrics satisfy common two-sided bounds
on $0\le s\le TL$, independently of $t$.

Apply Lemma~\ref{lem:comparison} to the two families and use
Proposition~\ref{prop:residual}. Uniformly for $0\le s\le TL$,
\begin{equation}\label{eq:global-sigma}
 \int_{X_t}\sigma(H_t(s),\widehat H_t(s))\,dA_t
 \le\eta_t+C_T\frac{(1+\log L)^2}{L}.
\end{equation}
Since $d^2\le\sigma$, this proves \eqref{eq:global-main}. The
integration is on the entire smooth compact fibre; no boundary
terms from truncated components have been omitted.

Fix $K_i\Subset U_i\Subset C_i\setminus\operatorname{Sing}X_0$.
The transported base geometry is uniformly smooth and elliptic on
$U_i$. For small $t$, the approximate residual there is
$O_T(L^{-2})$. The nonnegative scalar function
$u=\sigma(H_t,\widehat H_t)$ therefore satisfies
\[
 (\partial_s+\Delta_t)u\le C_TL^{-2}
\]
on that region. Choose finitely many interior coordinate balls
covering $K_i$, with doubled balls contained in $U_i$ and radii
at most $1/2$. All choices are independent of $t$.

For a backwards parabolic cylinder ending at $s_0\ge1$, the
function $u+C_TL^{-2}(s_0-s)$ is a nonnegative homogeneous
subsolution. The local scalar parabolic mean-value inequality
\cite[Theorem~5.1]{saloff-coste}, with exponent one, gives
\[
 \sup_{K_i(t)}u(s_0)
 \le C_{K_i,T}\left(
       \int_{s_0-1}^{s_0}\!\int_{X_t}u\,dA_t\,ds+L^{-2}\right).
\]
The ball radii and coefficient bounds determine the constant. Only
fixed compact charts are used, so no global Sobolev constant of the
degenerating fibre enters. The bound holds up to the upper time
endpoint by continuity. Substituting \eqref{eq:global-sigma} and
then using \eqref{eq:approx-thick-intro} and the triangle inequality
for $d$ proves \eqref{eq:local-main}.
\end{proof}

\begin{corollary}\label{cor:initial-components}
Suppose $cI\le H_t^{\mathrm{in}}\le CI$ and, on every fixed compact
component subset, the initial metric matrices converge in $L^2$ to
$H_i^0$. Then the initial-error hypothesis of Theorem~\ref{thm:main}
holds. In particular initial derivatives may be unbounded as $t\to0$.
\end{corollary}
\begin{proof}
On a common compact range of positive matrices, $\sigma$ is bounded
by a constant times squared Euclidean matrix difference. On each
fixed thick subset, $\widehat H_t(0)$ converges uniformly to $H_i^0$,
so the integral of $\sigma(H_t^{\mathrm{in}},\widehat H_t(0))$ there
tends to zero. The trace function is uniformly bounded everywhere.
The omitted node regions have area at most $C\rho_0^2$ for a fixed
small outer cutoff $\rho_0$, by Lemma~\ref{lem:plumbing}.
First send $t\to0$ for fixed $\rho_0$, then $\rho_0\to0$.
\end{proof}

The local estimate uses a fixed physical transient of length one;
any fixed positive transient can be obtained by choosing smaller
interior cylinders. On a fixed interval $\delta\le\tau\le T$ with
$\delta>0$, this transient is automatically passed for small $t$.
The global estimate applies from the initial time. It allows bounded
initial variation on the neck because its area tends to zero, but
does not imply pointwise approximation of that variation at time zero.

The finite equation is the leading evolution on the specified time
scale. The reconstructed metrics contain additional component and
endpoint corrections needed to make their accumulated residual small.
The argument does not identify the finite equation with an exact
invariant-manifold restriction of HYM at fixed $t$, and does not
replace finite $T$ by an unbounded slow-time interval.
